% Nicolás. S13--S15. Cómo lo probamos: datos reales, prototipo del relay,
% experimento de tamaño y simulación de la tormenta.
\section{Cómo lo probamos}
\begin{frame}{Paso 1: datos reales y un prototipo del relay}
  \centering
  \stat{70 min}{de BTC/USDT real}\hfill
  \stat{41\,113}{mensajes, \textbf{0} perdidos}\hfill
  \stat{9.8 / s}{mensajes de 5.8\,KB}\hfill
  \stat{275}{niveles por mensaje}\hfill
  \stat{41\,KB}{snapshot, 2\,000 niveles}\par
  \medskip
  % S13. Del dato real a la verificación del relay.
\begin{tikzpicture}[x=1cm,y=1cm]
  \path[use as bounding box] (0,1.1) rectangle (13.5,2.75);
  \tikzset{paso/.style={box,font=\scriptsize,text width=2.75cm,minimum height=13mm}}
  \node[paso] (p1) at (1.55,1.9) {\textbf{Grabación real}\\Binance Futures\\13 nov.\ 2026};
  \node[paso] (p2) at (5.0,1.9) {\textbf{Relay prototipo}\\secuencia y tombstone\\por nivel};
  \node[paso] (p3) at (8.45,1.9) {\textbf{46 reconexiones}\\al azar, cortes\\de distinto largo};
  \node[bueno,font=\scriptsize,text width=2.75cm,minimum height=13mm] (p4) at (11.9,1.9) {\textbf{46 / 46 correctas}\\checksum y 100 mejores\\niveles idénticos};
  \draw[flow] (p1) -- (p2);
  \draw[flow] (p2) -- (p3);
  \draw[flow,colReq] (p3) -- (p4);
\end{tikzpicture}

  \note{Nicolás. 10:40--11:30 (50 s). Para probarlo no inventamos el tráfico.
  Grabamos 70 minutos del stream público de profundidad de Binance Futures
  para BTC/USDT, el 13 de noviembre: 41 113 mensajes, sin perder ninguno.
  Llegan casi diez por segundo, y cada uno cambia en promedio 275 niveles. El
  snapshot tenía 2 000 niveles y pesaba 41 KB. Con esa grabación construimos
  un prototipo del relay, que aplica cada delta y guarda por nivel el último
  mensaje y las lápidas. Lo validamos con 46 reconexiones al azar: en las 46,
  el libro del cliente quedó idéntico al del relay, tanto en el checksum como
  en los 100 mejores niveles de cada lado.}
\end{frame}

\begin{frame}{Paso 2: ¿cuánto pesa recuperar a un cliente?}
  \begin{columns}[c]
    \begin{column}{0.6\textwidth}
      \centering% S14. Tamaño medio de una recuperación según la duración del corte, a ritmo
% normal. Datos: experimento de tamaño sobre la grabación (KB = 1024 bytes).
\begin{tikzpicture}
  \begin{axis}[
    width=8.6cm,height=5.0cm,
    xmode=log,ymode=log,
    xmin=0.8,xmax=75,ymin=8,ymax=40000,
    xtick={1,3,5,10,30,60},xticklabels={1,3,5,10,30,60},
    ytick={10,100,1000},yticklabels={{10\,KB},{100\,KB},{1\,MB}},
    xlabel={segundos desconectado},ylabel={tamaño de la recuperación},
    tick label style={font=\scriptsize},label style={font=\scriptsize},
    axis line style={colNeutro},
    legend style={font=\scriptsize,draw=none,fill=none,at={(0.02,0.98)},anchor=north west},
    legend cell align=left]
    \addplot[colNeutro,thick,mark=*,mark size=1.4pt] coordinates {
      (1,44.1) (3,147.5) (5,252.5) (10,540.6) (30,1313.5) (60,3433.1)};
    \addlegendentry{Estilo Ably: todos los mensajes}
    \addplot[colWake,thick,dashed] coordinates {
      (1,41.1) (3,41.1) (5,41.1) (10,41.1) (30,41.1) (60,41.1)};
    \addlegendentry{Binance: snapshot completo}
    \addplot[colReq,line width=1.6pt,mark=*,mark size=1.6pt] coordinates {
      (1,15.4) (3,31.2) (5,38.0) (10,41.1) (30,41.1) (60,41.1)};
    \addlegendentry{RelayBook: niveles cambiados}
  \end{axis}
\end{tikzpicture}

    \end{column}
    \begin{column}{0.38\textwidth}\small
      Medido sobre la grabación, a ritmo normal:\par\medskip
      \textbf{1\,s fuera:} 15\,KB vs 41\,KB del snapshot, \textbf{63\% menos}.\par\medskip
      \textbf{Frente a Ably:} entre \textbf{3 y 84 veces} menos; Ably llega a 3.4\,MB tras 60\,s.\par\medskip
      \textbf{Desde 5--10\,s:} cambiaron tantos niveles que el relay manda el snapshot.
    \end{column}
  \end{columns}
  \takeaway{Reanudar por estado gana en cortes breves y siempre gana a reenviar mensajes.}
  \note{Nicolás. 11:30--12:30 (60 s). Primer experimento: ¿cuánto pesa recuperar a
  un cliente? Sobre la grabación, a ritmo normal, desconectamos clientes en
  momentos al azar, por 1 a 60 segundos, y medimos cuánto hay que mandarle con
  cada diseño. El gráfico está en escala logarítmica. La línea roja punteada
  es el snapshot de Binance: siempre 41 KB. La gris es el estilo Ably: crece
  porque reenvía cada mensaje, y llega a 3.4 MB tras un minuto. La verde es
  nuestro resume por estado: tras un segundo fuera, 15 KB, un 63 \% menos que
  el snapshot. Entre 5 y 10 segundos ya cambiaron tantos niveles que el relay
  manda el snapshot, así que nunca pagamos más que el libro. Frente a Ably,
  entre 3 y 84 veces menos. Funciona porque muchos mensajes vuelven a tocar
  los mismos niveles: en 10 segundos llegan unas 27 000 actualizaciones, pero
  solo a unos 2 800 niveles distintos.}
\end{frame}

\begin{frame}{Paso 3: simulamos la tormenta con el tráfico real}
  \centering% S15. Escenario de la simulación de la tormenta y qué se barrió.
\begin{tikzpicture}[x=1cm,y=1cm]
  \path[use as bounding box] (0,0.15) rectangle (13.5,4.7);
  \node[font=\scriptsize] at (6.75,4.5) {Mismo tráfico para \textbf{Binance}, \textbf{estilo Ably} (ventana de 2 min) y \textbf{RelayBook}};
  % Línea de tiempo: 1 cm = 1 min 40 s
  \draw[draw=colNeutro,fill=colNeutro!8] (0.75,3.6) rectangle (3.75,4.1);
  \draw[draw=colWake,fill=colWake!18] (3.75,3.6) rectangle (6.75,4.1);
  \draw[draw=colNeutro,fill=colNeutro!8] (6.75,3.6) rectangle (12.75,4.1);
  \node[font=\scriptsize] at (2.25,3.85) {5 min a 1$\times$};
  \node[font=\scriptsize\bfseries,text=colWake] at (5.25,3.85) {5 min a 10$\times$, pico};
  \node[font=\scriptsize] at (9.75,3.85) {10 min a 1$\times$, recuperación};
  \draw[flow,colWake] (3.75,3.0) -- (3.75,3.55);
  \node[font=\scriptsize,text=colWake,anchor=west] at (3.85,3.15) {el 30\% de los clientes se desconecta y vuelve a los 1--5\,s};
  % Parámetros
  \tikzset{par/.style={box,font=\scriptsize,text width=2.9cm}}
  \node[par] at (1.65,2.2) {\textbf{2\,000 clientes}\\3 relays};
  \node[par] at (5.05,2.2) {\textbf{70\%} rápidos a 2\,MB/s\\\textbf{30\%} lentos a 80\,KB/s};
  \node[par] at (8.45,2.2) {\textbf{Cola de 1\,MB}\\Binance y Ably cortan\\conflation en RelayBook\\desde 256\,KB};
  \node[par] at (11.85,2.2) {\textbf{Retries}\\reconecta tras 5\,s sin datos\\reintenta tras 10\,s};
  % Barrido
  \node[font=\scriptsize\bfseries,text=colAcento] at (6.75,1.35) {Parámetros desconocidos de Binance probados con varios valores, 6 $\times$ 3 = 18 corridas};
  \node[acento,font=\scriptsize,text width=5.9cm] at (3.35,0.6) {\textbf{Capacidad del servidor}\\de 1.5$\times$ a 15$\times$ la carga normal};
  \node[acento,font=\scriptsize,text width=5.9cm] at (10.15,0.6) {\textbf{Costo de generar un snapshot}\\1$\times$, 10$\times$ y 50$\times$ lo que cuesta enviarlo};
\end{tikzpicture}

  \takeaway{Capacidad y costo de Binance son desconocidos: probamos varios valores.}
  \note{Nicolás. 12:30--13:25 (55 s). Segundo experimento, el central: simular la
  tormenta. El tiempo avanza en pasos de 100 milisegundos y cada servidor
  reparte su capacidad entre sus conexiones. Reproducimos la grabación en tres
  fases: cinco minutos normales, cinco a diez veces y diez normales, para ver
  si el sistema se recupera solo. Al inicio del pico se desconecta el 30 \% de
  los 2 000 clientes, y el 30 \% tiene internet lento. Si una conexión acumula
  1 MB pendiente, Binance y Ably la cortan; RelayBook aplica conflation desde
  los 256 KB. El estilo Ably corre sobre los mismos relays, así que solo
  cambia cómo reanuda. Dos cosas no las conocemos de Binance: cuánta capacidad
  le sobra y cuánto le cuesta generar un snapshot. Así que probamos varios
  valores: seis capacidades por tres costos, 18 corridas.}
\end{frame}
